% !PRINTED 410
\weberpage{410}
%<<CONT
so folgt, dass $r^\mu = 1$ sein muss. Es kann also in der Reihe (6) kein früher
dagewesenes Glied wiederkehren, ehe das erste Glied 1 zum zweiten Male vorkommt, und
wenn $\mu$ die kleinste positive Zahl ist, für die $r^\mu = 1$ ist, so sind die Zahlen:

\weq{7}{1, r, r^2 \ldots r^{\mu - 1}}
alle von einander verschieden.

\gesperrt{Es muss dann $\mu$ ein Theiler von $n$ sein.} Denn durch Division lassen sich die ganzen
Zahlen $h, \mu'$ so bestimmen, dass

\[
n = h\mu + \mu'; \quad 0 \leqq \mu' < \mu.
\]
Dann ist aber auch, wie aus

\[
r^n = r^{h\mu} r^{\mu'}
\]
hervorgeht, $r^{\mu'} = 1$, d. h. da $\mu$ die kleinste positive Zahl sein soll, für die
$r^\mu = 1$ ist, $\mu' = 0$, oder $n$ durch $\mu$ theilbar.

\gesperrt{Es ist also $r$ zugleich $\mu^{\mathrm{te}}$ Einheitswurzel, aber nicht Einheitswurzel
von noch niedrigerem Grade.}

\gesperrt{Man nennt die Zahl $r$ eine primitive $n^{\mathrm{te}}$ Einheitswurzel, wenn
sie nicht zugleich Einheitswurzel eines niedrigeren Grades ist.}

Aus dieser Definition folgt, dass die Zahlen der Reihe (6), wenn $r$ eine primitive
$n^{\mathrm{te}}$ Einheitswurzel ist, alle von einander verschieden sind, und dass
sämmtliche $n^{\mathrm{ten}}$ Einheitswurzeln darunter enthalten sind, dass sie aber nur
einen Theil der $n^{\mathrm{ten}}$ Einheitswurzeln ausmachen, wenn $r$ eine imprimitive
$n^{\mathrm{te}}$ Einheitswurzel ist.

Jede Einheitswurzel, deren Grad ein von $n$ verschiedener Theiler von $n$ ist, ist
zugleich imprimitive $n^{\mathrm{te}}$ Einheitswurzel.

\gesperrt{Ist $r$ zugleich $n^{\mathrm{te}}$ und $m^{\mathrm{te}}$ Einheitswurzel, so ist es auch
$\mu^{\mathrm{te}}$ Einheitswurzel, wenn $\mu$ der gr\"osste gemeinschaftliche Theiler von
$n$ und $m$ ist.}

Denn nach \S.\,118 kann man die ganzen Zahlen $x, y$ so bestimmen, dass

\[
m x + n y = \mu
\]
wird, und folglich ist

\[
r^\mu = r^{mx} r^{ny} = 1.
\]
Dem entspricht der andere Satz:

\gesperrt{Sind $r_1, r_2 \ldots$ Einheitswurzeln der Grade $n_1, n_2 \ldots$, so sind sie alle
zugleich Einheitswurzeln des Grades $m$, wenn $m$ irgend ein gemeinschaftliches
Vielfaches von $n_1, n_2 \ldots$ bedeutet.}
